\documentclass[11pt]{amsart}
\usepackage{amsmath,amssymb,amsthm,mathtools}
\usepackage[margin=1in]{geometry}

\newtheorem{question}{Question}
\theoremstyle{definition}
\newtheorem{fact}{Fact}

\newcommand{\PSL}{\mathrm{PSL}}
\newcommand{\PGL}{\mathrm{PGL}}
\newcommand{\SL}{\mathrm{SL}}
\newcommand{\GL}{\mathrm{GL}}
\newcommand{\R}{\mathbb{R}}
\newcommand{\Z}{\mathbb{Z}}
\newcommand{\C}{\mathbb{C}}
\newcommand{\Phin}[1]{\Phi_{#1}}

\title{A rigidity step in the cyclotomic argument for $q$-deformed Pell automorphs}
\date{21 July 2026}

\begin{document}
\maketitle

\noindent
This note isolates one step of a proposed argument. Everything else in the
argument is either classical or already proved in the current draft. The
question at the end is the only thing being asked.

\section*{Setting}

For a nonsquare $d$ with fundamental Pell solution $(r,s)$, let
$\Lambda_q=\begin{psmallmatrix}A&B\\C&D\end{psmallmatrix}$ be the $q$-deformed
automorph, entries in $\Z[q,q^{-1}]$, specialising to
$\begin{psmallmatrix}r&ds\\s&r\end{psmallmatrix}$ at $q=1$. The statement of
interest, called Lemma~L in the draft, is
\[
  \Phin{n}\mid C \quad\Longrightarrow\quad A\equiv D \pmod{\Phin{n}} .
\]
It is proved for $n\in\{2,3,4,5\}$ and for $n$ an odd prime. The open range is
everything else, and the prime powers $p^{a}$ with $a\ge 2$ are the case the
draft needs.

\begin{fact}[from the draft]
The assignment $R\mapsto R_q=\begin{psmallmatrix}q&1\\0&1\end{psmallmatrix}$,
$S\mapsto S_q=\begin{psmallmatrix}0&-1/q\\1&0\end{psmallmatrix}$ is
projectively multiplicative on $\PSL(2,\Z)$. At $q=\zeta$ a primitive $n$th
root of unity, $R_\zeta^{\,n}=I$ and the image of $R$ has projective order
exactly $n$.
\end{fact}

\begin{fact}[Coxeter and Moser]
With $a:=S$ and $b:=(SR)^{-1}$ one has $\PSL(2,\Z)=\Z/2*\Z/3$, and imposing
$R^{n}=1$ gives the von Dyck group
$D(2,3,n)=\langle a,b \mid a^{2}=b^{3}=(ab)^{n}=1\rangle$.
\end{fact}

So the specialised projective representation
$\bar\rho:\PSL(2,\Z)\to\PGL(2,\C)$ factors through $D(2,3,n)$.

\section*{Why the existing proof stops}

For $n\le 5$ the image is finite and the argument runs through $\PSL(2,\Z/n)$,
using the identification of the normal closure $\langle\langle R^{n}\rangle\rangle$
with $\pm\Gamma(n)$. For $n\ge 6$ that identification is false: the normal
closure is a noncongruence subgroup of infinite index. The proposed repair is
that the finite quotient was never the load-bearing object, and that what is
actually needed is the injectivity of the induced map
$\pi: D(2,3,n)\to\PGL(2,\C)$.

\section*{The step in question}

Fix $n\ge 7$ and $\zeta=e^{2\pi i/n}$. Take the $\SL(2,\C)$ lifts
\[
  \widetilde S=\zeta^{1/2}S_\zeta,\qquad
  \widetilde R=\zeta^{-1/2}R_\zeta,\qquad
  \widetilde B=(\widetilde S\widetilde R)^{-1},
\]
all of determinant $1$. On the generating pair $(a,b)$ the trace triple is
\[
  \bigl(\operatorname{tr}\widetilde S,\ \operatorname{tr}\widetilde B,\
   \operatorname{tr}\widetilde S\widetilde B\bigr)
  =\bigl(0,\ 1,\ 2\cos(\pi/n)\bigr),
\]
since $\widetilde S\widetilde R=\begin{psmallmatrix}0&-1/\zeta\\ \zeta&1\end{psmallmatrix}$
has trace $1$, and $\widetilde S\widetilde B$ is conjugate to $\widetilde R^{-1}$.
The reducibility invariant is
\[
  x^{2}+y^{2}+z^{2}-xyz-4 \;=\; 4\cos^{2}(\pi/n)-3 ,
\]
which vanishes only at $n=6$. So for $n\ne 6$ the representation is irreducible.

\begin{fact}[Poincar\'e]
For $n\ge 7$ the hyperbolic $(\pi/2,\pi/3,\pi/n)$ triangle generates a
cocompact Fuchsian group whose orientation-preserving subgroup is isomorphic to
$D(2,3,n)$ via the presentation map, so this projective representation is
faithful. Its $\SL(2,\R)$ lifts have traces $0$, $\pm 1$, $\pm 2\cos(\pi/n)$.
\end{fact}

\begin{fact}[Fricke]
Two irreducible $\SL(2,\C)$ representations of a free group of rank two with
equal trace triples are conjugate in $\GL(2,\C)$, and the lift sign changes
realise $(x,y,z)\mapsto(-x,y,-z)$ and $(x,-y,-z)$.
\end{fact}

The proposed conclusion is that the two representations, having the same trace
triple and both being irreducible, are conjugate; that conjugacy descends
projectively; and that faithfulness of the Fuchsian one therefore forces $\pi$
to be injective, giving
$\ker\bar\rho=\langle\langle R^{n}\rangle\rangle$. From there the rest of the
argument is arithmetic and is not in doubt.

\begin{question}
Is that inference valid as stated, and in particular:
\begin{enumerate}
\item Fricke rigidity is a statement about representations of a free group of
rank two, whereas both representations here are representations of
$\Z/2 * \Z/3$ pulled back along the natural surjection. Does conjugacy of the
pullbacks give conjugacy of the representations of $D(2,3,n)$ that is needed
for the kernel comparison?
\item The order-$n$ elliptic generator of a Fuchsian $(2,3,n)$ group has trace
$\pm 2\cos(\pi k/n)$ for some $k$ coprime to $n$, and only $k=1$ corresponds to
the discrete faithful realisation. The specialisation at $\zeta=e^{2\pi i/n}$
produces $k=1$, but at $\zeta^{k}$ it produces the triple
$(0,1,2\cos(\pi k/n))$, which should correspond to a representation that is
neither discrete nor faithful. Is it legitimate to prove faithfulness at the
distinguished root only, and does anything need saying about the other
conjugates?
\item A related step identifies the upper-triangular elements of the image with
the stabiliser of a point, concluding it is cyclic of order exactly $n$ because
finite subgroups of a cocompact triangle group are conjugate into vertex
stabilisers. Is that appeal standard enough to cite, or does it need an
argument?
\end{enumerate}
\end{question}

\section*{What is at stake}

If the step holds, Lemma~L becomes unconditional at every $n\ge 7$ with
$n\ne 6$, prime powers and composites alike, and the downward-closure
hypothesis currently assumed in the draft is not needed for it at all. Two
further consequences follow: $\Phin{n}\mid B$ whenever $\Phin{n}\mid C$, which
extends the proved fragment of the scalar-congruence conjecture from
$\{2,3,4,5\}$ to all indices except $6$; and the conditional reduction at odd
prime powers inverts, so that the boundary congruence
$p\mid\gamma_{p^{a-1}}$ becomes a corollary rather than an open input.

The index $n=6$ is not reachable by this argument, the trace triple falling
exactly on the reducible locus, and it is also the one index where the
alternative $\pi$-adic route degenerates, since $\Phin{6}(1)=1$.

Numerically the three conclusions have been checked with exact arithmetic on
$213$ instances, and on a further $116$ instances at $n=8,16,32$ by an
independent implementation, with no violations. That is consistent with the
argument but does not test the step above.

\end{document}
